\item Un submarino está a $300\text{ m}$ de distancia de la orilla de un lago de agua dulce y $100\text{ m}$ por debajo de la superficie del agua. Desde el submarino se envía un rayo láser de tal manera que el haz incide en la superficie del agua a $210\text{ m}$ de la costa. Hay un edificio en la orilla y el haz láser incide en un blanco en la parte superior del edificio.
\begin{enumerate}
    \item Dibuje un diagrama de la situación, identificando los dos triángulos que son importantes en la búsqueda de la solución.
    \item Encuentre el ángulo de incidencia del rayo al golpear la interfaz agua-aire.
    \item Encuentre el ángulo de refracción.
    \item ¿Qué ángulo forma el rayo refractado con la horizontal?
    \item Determine la altura del objetivo sobre la superficie del agua (nivel del lago).
\end{enumerate}